\chapter{Image-Guided Drainage}

\label{ch:ch78_drainage}

Image-guided drainage is a minimally invasive procedure in which a catheter is placed into an abnormal fluid collection under imaging guidance. It provides source control for infection, relieves obstruction, and avoids open surgery in many patients.

---

\section{1. Patient Perspective: What a Drain Insertion Involves}

If you have a collection of infected fluid (abscess) or a blocked bile duct, your doctor may suggest inserting a \textbf{drain} using imaging guidance.

\begin{itemize}
  \item The procedure is performed under local anaesthetic, sometimes with sedation.
  \item Using CT or ultrasound, the doctor locates the fluid collection and plans a safe path for the drain.
  \item A small skin incision is made. A needle is advanced into the collection, and a thin plastic catheter is threaded over a guide wire into position (Seldinger technique).
  \item Fluid is drained, and samples are sent for laboratory analysis.
  \item The catheter remains in place and is connected to a drainage bag. It will be removed once output decreases and the collection resolves.
  \item The procedure takes 30--60 minutes. Most patients go home the same day or after an overnight stay.
\end{itemize}

---

\section{2. Indications}

\begin{itemize}
  \item \textbf{Abscess drainage}: Intra-abdominal abscess (liver, splenic, diverticular, appendiceal), pelvic abscess, or mediastinal abscess.
  \item \textbf{Biliary drainage} (percutaneous transhepatic biliary drainage, PTBD): Obstructive jaundice due to malignant or benign biliary stricture when endoscopic drainage has failed or is not feasible.
  \item \textbf{Nephrostomy}: Urinary obstruction (hydronephrosis) when ureteric stenting has failed or is contraindicated.
  \item \textbf{Pleural drainage}: Empyema, complicated parapneumonic effusion, or chylothorax.
  \item \textbf{Other}: Lymphocele drainage post-transplant, pseudocyst drainage, and infected necrotic collections.
\end{itemize}

\section{3. Catheter Types and Selection}

\begin{table}[h]
\centering
\begin{tabular}{l l l}
\hline
\textbf{Catheter Type} & \textbf{Size} & \textbf{Common Use} \\
\hline
All-purpose drainage (APD) & 12--16 Fr & Larger abscesses, infected collections \\
Malecot catheter & 14--24 Fr & Nephrostomy, large abscesses \\
Pigtail catheter & 8--12 Fr & Small abscesses, simple collections, biliary \\
Pleural catheter & 24--32 Fr & Empyema, large pleural effusions \\
\hline
\end{tabular}
\caption{Catheter Selection for Image-Guided Drainage}
\end{table}

\begin{pearlbox}[Pigtail vs Large-Bore]
Use a \textbf{pigtail catheter} for simple, low-viscosity collections (serous fluid, thin pus). Use a \textbf{large-bore catheter} (APD or Malecot) for thick, loculated, or haemorrhagic collections that require larger calibre drainage.
\end{pearlbox}

---

\section{4. The Seldinger Technique}

The Seldinger technique is the standard method for catheter placement:

\begin{enumerate}
  \item \textbf{Access}: A needle is advanced into the collection under imaging guidance.
  \item \textbf{Wire}: A flexible guide wire is passed through the needle into the collection.
  \item \textbf{Dilatation}: The tract is dilated over the wire using serial dilators.
  \item \textbf{Catheter}: The drainage catheter is advanced over the wire into the collection.
  \item \textbf{Wire removal}: The guide wire is removed, leaving the catheter in position.
  \item \textbf{Secure and connect}: The catheter is secured to the skin and connected to a drainage bag.
\end{enumerate}

---

\section{5. Drainage vs Surgery: When to Recommend Each}

\begin{itemize}
  \item \textbf{Percutaneous drainage is preferred} when the collection is well-defined, unilocular, and accessible; the patient is haemodynamically stable; and surgical risk is high.
  \item \textbf{Surgery is preferred} when there is diffuse peritonitis, contained perforation requiring resection, failed percutaneous drainage after 48--72 hours, or multiloculated collections that cannot be adequately drained percutaneously.
  \item \textbf{Multidisciplinary discussion} is essential for complex cases.
\end{itemize}

\section{6. Catheter Management and Follow-Up}

\begin{itemize}
  \item \textbf{Daily assessment}: Catheter patency, output volume and character, skin site, and patient clinical status.
  \item \textbf{Flush protocol}: Gentle flush with 5--10 mL saline 1--2 times daily (varies by institution) to maintain patency.
  \item \textbf{Follow-up imaging}: Repeat CT or ultrasound to assess collection size. If the collection is not reducing, consider catheter repositioning, upsizing, or additional drainage.
  \item \textbf{Removal criteria}: Output $<$ 20--30 mL/day, clinical improvement (afebrile, normalising WCC), and follow-up imaging showing resolution of the collection.
\end{itemize}

\begin{pearlbox}[Clinical Pearl]
A drain that stops draining does not necessarily mean the collection has resolved. Always confirm with follow-up imaging before removing the catheter.
\end{pearlbox}

---

\section{7. Structured Reporting Example}

\begin{testbox}[Indications:]
55-year-old diabetic male with persistent fever despite 5 days of IV antibiotics. CT abdomen shows a 6 cm rim-enhancing fluid collection in the right lobe of the liver. Percutaneous drainage recommended.
\end{testbox}
\begin{refbox}[Findings:]
\textbf{Pre-procedure CT}: A 6.2 $\times$ 4.8 cm rim-enhancing, hypodense collection in segments VII/VIII of the liver, with surrounding oedema. No vascular involvement. The intercostal space and lung base are clear. Safe transhepatic route identified.\\
\textbf{Procedure}: Under CT guidance and local anaesthesia, an 8 Fr pigtail catheter was placed via a transhepatic approach using the Seldinger technique. Approximately 85 mL of thick, purulent fluid was aspirated. The catheter was secured and connected to a drainage bag.\\
\textbf{Post-procedure CT}: The catheter tip is well-positioned within the residual collection. No pneumothorax, haemoperitoneum, or new free fluid.
\end{refbox}
\begin{pearlbox}[Impression:]
1. Successful CT-guided percutaneous drainage of a hepatic abscess (6.2 cm).
2. Approximately 85 mL of pus aspirated; specimen sent for culture and sensitivity.
3. Catheter left in situ for ongoing drainage. Follow-up imaging in 48 hours.
\end{pearlbox}

\clearpage
